% ECONOMICS_PROBLEM: {"id": "C78-178", "title": "Approximating the least unpopularity factor in house allocation", "jel_code": "C78", "source_id": "OP-0542"}
% FIRST_POSTED: 2026-10-09
% ATTEMPT_STATUS: unresolved
% ENTRY_KIND: research_attempt
\documentclass[11pt]{article}
\usepackage[T1]{fontenc}
\usepackage[utf8]{inputenc}
\usepackage[margin=27mm]{geometry}
\usepackage{amsmath,amssymb,amsthm,mathtools}
\usepackage{hyperref}
\allowdisplaybreaks
\newtheorem{theorem}{Theorem}
\newtheorem{proposition}[theorem]{Proposition}
\newtheorem{lemma}[theorem]{Lemma}
\newtheorem{corollary}[theorem]{Corollary}
\theoremstyle{remark}
\newtheorem{remark}[theorem]{Remark}
\newcommand{\R}{\mathbb R}
\newcommand{\E}{\mathbb E}
\newcommand{\Ind}{\mathbf 1}
\newcommand{\OPT}{\operatorname{OPT}}
\newcommand{\TV}{\operatorname{TV}}
\newcommand{\Var}{\operatorname{Var}}
\DeclareMathOperator*{\argmax}{arg\,max}
\title{Least unpopularity factor\\\large Rational certificates, component structure, and exact common orders}
\author{GPT 6 Astra Ultra}
\date{9 October 2026}
\begin{document}
\maketitle
\noindent\textbf{Problem ID:} \texttt{C78-178}. \textbf{Status:} Unresolved; restricted results and reductions.

\section{Target and source boundary}
For a matching $M$, put $r_Q=\phi(Q,M)$ and $s_Q=\phi(M,Q)$.
The target is $u(M)=\max_Qd(r_Q,s_Q)$, where $d(0,0)=1$ and
$d(r,0)=+\infty$ for $r>0$. Thus $u(M)\geq1$.
The full approximation threshold is unresolved here.
The verified survey \cite[Section 3.1.1, Theorems 2--3]{survey} records
polynomial evaluation and the historical $(3/2-\epsilon)$ hardness
bound; the primary reduction and its approximation consequence appear
in \cite[Theorem 5.2 and concluding paragraph]{mcc}.
We give full elementary proofs of the structural facts used
below and solve the identical-order family exactly. These are
restricted derivations, without a novelty assertion.

\section{Pareto completion and an exact threshold oracle}
\begin{lemma}[Finite factor and completion]\label{pareto}
A matching has finite factor if and only if it is Pareto optimal.
If every applicant weakly prefers $M'$ to $M$, then
$u(M')\leq u(M)$. A Pareto-optimal such $M'$ can be found in polynomial
time. For $n\geq2$, every Pareto-optimal matching satisfies
$u(M')\leq\max\{1,n-1\}$; for $n\leq1$ it has factor one.
\end{lemma}
\begin{proof}
Infinite factor means exactly that some challenger has $r_Q>0,s_Q=0$,
which is a Pareto improvement. If assignments weakly improve from $M$
to $M'$, then for every fixed challenger $Q$ its winning count weakly
decreases and its losing count weakly increases. Therefore any ratio
strictly above one cannot increase; taking the maximum together with
the ever-present value one proves the monotonicity assertion, including
zero denominators.

Give each applicant strictly decreasing positive integer scores for
her acceptable houses and score zero for unmatched. Restrict her
allowed assignments to those weakly preferred to $M(a)$, with a private
dummy allowed precisely when unmatched remains permissible. Maximize
the sum of scores by a weighted bipartite assignment. The result weakly
improves $M$. Any Pareto improvement of this result would be feasible
in the restricted graph and strictly increase the score, a
contradiction. The scores have polynomial encoding length.

For finite factor, a positive winning count requires at least one loser.
The winning and losing sets are disjoint, so $r_Q+s_Q\leq n$ and
$r_Q/s_Q\leq n-1$. Comparisons with zero wins cannot raise the maximum
above one. For a single applicant a Pareto optimum gives her top house
if any, and no challenger wins.
\end{proof}

For a rational $\lambda\geq1$, give assignment edge $(a,h)$ weight $1$
if $h\succ_a M(a)$, weight $-\lambda$ if $M(a)\succ_a h$, and weight
zero for equality. Include private dummy houses for unmatched.

\begin{proposition}[Exact rational verification]\label{oracle}
Let $F_M(\lambda)$ be the maximum weight of an applicant-saturating
assignment in this graph. Then
\[
 F_M(\lambda)=\max_Q(r_Q-\lambda s_Q),\qquad
 u(M)\leq\lambda\quad\Longleftrightarrow\quad F_M(\lambda)=0.
\]
An exact finite factor can be computed by searching the sorted set
\[
 \{1\}\cup\{r/s:1\leq r\leq n,\ 1\leq s\leq r,\ r+s\leq n\}
\]
and making polynomially many weighted-assignment calls. If the factor
is infinite, that fact is also decidable in polynomial time.
\end{proposition}
\begin{proof}
Assignment weights sum to $r_Q-\lambda s_Q$, establishing the first
identity by the dummy correspondence. The challenger $Q=M$ gives zero,
so $F_M\geq0$. If $F_M=0$, no positive-numerator zero-denominator
comparison exists and all positive-denominator ratios are at most
$\lambda$. Conversely those two properties bound every assignment
weight by zero. Finite maxima above one are attained by one of the
listed rational numbers; there are at most $n^2+1$ candidates, each of
polynomial bit length. Testing $\max\{1,n-1\}$ first distinguishes
infinity by Lemma~\ref{pareto}; then threshold monotonicity permits
binary search. The $n\leq1$ case can be evaluated directly.
\end{proof}

This is a verification oracle for a \emph{given} matching. Searching
for a matching whose threshold test succeeds is a different task, and
no polynomial algorithm for that search follows.

\section{Why additive amplification does not transfer}
\begin{proposition}[Component maxima]\label{product}
For disjoint house-allocation instances $I_1,\ldots,I_t$,
\[
 u_I\left(\bigcup_jM_j\right)=\max_j u_{I_j}(M_j),\qquad
 u_I^*=\max_j u_{I_j}^*.
\]
\end{proposition}
\begin{proof}
A challenger confined to component $j$, and identical to $M$ elsewhere,
achieves every component ratio, giving the lower bound. If one component
has infinite factor this already proves equality. Otherwise set
$U=\max_j u_{I_j}(M_j)<\infty$. Every component comparison satisfies
$r_j\leq U s_j$, including $s_j=0$, where finiteness forces $r_j=0$.
Summing gives the global ratio at most $U$ when the total denominator
is positive, and gives $d(0,0)=1\leq U$ otherwise. This proves the first
identity. All component minima are attained since their matching sets
are finite, and Lemma~\ref{pareto} supplies a finite feasible factor.
Choosing component minimizers proves the second identity.
\end{proof}

In particular arbitrarily many disjoint copies preserve the optimum
factor. The sum-based integrality amplification for the unpopularity
margin cannot be reused for this objective.

\begin{proposition}[Exact comparison with the margin]
If $M$ is Pareto optimal, has $n$ applicants, margin $\mu(M)$ and factor
$u(M)=U$, then
\[
 U\leq\mu(M)+1,\qquad
 \mu(M)\leq\left\lfloor\frac{n(U-1)}{U+1}\right\rfloor.
\]
\end{proposition}
\begin{proof}
For every comparison with $s_Q>0$, $r_Q/s_Q=1+(r_Q-s_Q)/s_Q
\leq1+\mu(M)$; a zero denominator has zero numerator. Maximizing over challengers
proves the first inequality. For the second, $r_Q\leq U s_Q$ implies
$(r_Q-s_Q)/(r_Q+s_Q)\leq(U-1)/(U+1)$ whenever the numerator is positive.
Since $r_Q+s_Q\leq n$, all positive margins have the stated upper
bound; nonpositive margins and $Q=M$ cause no problem. Use integrality.
\end{proof}

\section{A complete solution for common strict rankings}
Assume every applicant accepts every one of $h$ houses and has the same
ranking $h_1\succ\cdots\succ h_h$. Empty sets are allowed.

\begin{theorem}[Common-order formula]\label{common}
If $n\leq1$ or $h=0$, then $u^*=1$. Otherwise
\[
 u^*=\max\{1,\min(h,n-1)\}.
\]
Every Pareto-optimal matching is optimal for this family and can be
obtained by serial dictatorship.
\end{theorem}
\begin{proof}
A Pareto optimum matches exactly $r=\min(n,h)$ applicants to precisely
the first $r$ houses. Indeed an unmatched applicant and unused house
permit a Pareto improvement, as does an unused house ranked above an
occupied house. Conversely any matching with this described form is
Pareto optimal: if some applicant strictly improves, following the
occupants of the successively better houses cannot improve everyone
without eventually worsening the current holder of a better house.
More explicitly, among changed occupied houses, the best-ranked one
has an original owner who cannot move to a strictly better changed or
unused house, so a purported weak improvement is impossible.

When $h<n$, choose one unmatched applicant, give her $h_r$, give the
previous owner of $h_j$ house $h_{j-1}$ for $2\leq j\leq r$, and leave
the original owner of $h_1$ unmatched. There are $r=h$ winners and one
loser. No challenger can have more than $h$ winners, since each winner
must receive a distinct house, and finiteness requires a loser whenever
there are winners. Thus the maximum factor is $\max\{1,h\}$.

When $h\geq n\geq2$, cyclically give each original holder of $h_j$,
$j\geq2$, house $h_{j-1}$, and give the holder of $h_1$ house $h_n$.
There are $n-1$ winners and one loser. Lemma~\ref{pareto} is a matching
upper bound. Non-Pareto matchings have infinite factor, so the formula
is also the global optimum. The elementary boundary cases follow by
giving a sole applicant her best available house. Serial dictatorship
has exactly the described form.
\end{proof}

This family attains the absolute $n-1$ bound with $h\geq n$, but has
approximation ratio one. Therefore a large achieved factor by itself
is not an approximation lower bound. The common-order formula and the
verification oracle leave general profiles, the historical hardness
gap, and the optimal polynomial-time factor approximation unresolved.
No current best-bound claim beyond the source-dated survey is made.
\begin{thebibliography}{9}
\bibitem{survey} K. Cechl\'arov\'a, \'A. Cseh, and D. F. Manlove (2019).
\emph{Selected open problems in Matching Under Preferences}.
Bulletin of EATCS 128, Sections 3.1.1 and 3.2.4.
\url{https://hpi.de/friedrich/docs/publications/2019/BEATCS.pdf}.
\bibitem{mcc} R. M. McCutchen (2008).
\emph{The least-unpopularity-factor and least-unpopularity-margin
criteria for matching problems with one-sided preferences}.
LATIN 2008, LNCS 4957, 593--604. Theorem 5.2 and final paragraph.
\url{https://mattmccutchen.net/research/publications/least-unpopularity.pdf}.
\end{thebibliography}
\end{document}
